\documentclass[11pt]{article}
\usepackage{vinotheca}

\begin{document}

\vinothecaTitle{Region Resonances}{Data Appendix}{Source Data, Tables, and Reference Materials}
\vinothecaAuthor{Jure Skarabot}{May 2026}
\vinothecaCompanions{Summary \textperiodcentered{} Technical Appendix \textperiodcentered{} Methods Primer \textperiodcentered{} Data Appendix \textperiodcentered{} Region Reference (Identity)}

\section*{Table of Contents}
\begin{itemize}[leftmargin=2em,itemsep=0.1em,topsep=0.1em]
\item D.1 Region Roster (59 regions, with metaphors)
\item D.2 Embedding Model Specification
\item D.3 Calibration Test Set and Operating Range
\item D.4 Display Transformation Constants
\item D.5 Distribution of Top-$k$ Cosine Similarities
\end{itemize}

\section*{D.1 Region Roster}
The 59 regions analysed by Region Resonances, listed with their canonical soul metaphors. The corpus is identical to that used by Region Affinities; the metaphors are drawn from the canonical Layer 1 identity descriptions and serve as the first-token anchor of each region's embedded passage.

\begin{center}
\renewcommand{\arraystretch}{1.05}
\small
\begin{tabular}{@{}r l l l@{}}
\toprule
\# & Region & Country & Metaphor \\
\midrule
1  & Kamptal              & Austria    & Discipline \\
2  & Steiermark           & Austria    & Clarity \\
3  & Wachau               & Austria    & Monumentality \\
4  & Wagram               & Austria    & Earth \\
5  & Dalmatian Coast      & Croatia    & Tranquility \\
6  & Alsace               & France     & Duality \\
7  & Beaujolais           & France     & Joy \\
8  & Bordeaux             & France     & Business \\
9  & Burgundy             & France     & Devotion \\
10 & Champagne            & France     & Society \\
11 & Ch\^ateauneuf-du-Pape & France    & Family \\
12 & Jura                 & France     & Eccentricity \\
13 & Loire                & France     & Sentimentality \\
14 & Northern Rh\^one     & France     & Solitude \\
15 & Provence             & France     & Pleasure \\
16 & Baden                & Germany    & Warmth \\
17 & Mosel                & Germany    & Poetry \\
18 & Nahe                 & Germany    & Subtlety \\
19 & Pfalz                & Germany    & Generosity \\
20 & Rheingau             & Germany    & Nobility \\
21 & Macedonia            & Greece     & Austerity \\
22 & Santorini            & Greece     & Survival \\
23 & Tokaj                & Hungary    & Melancholy \\
24 & Alto Adige           & Italy      & Precision \\
25 & Campania             & Italy      & Memory \\
26 & Etna                 & Italy      & Awakening \\
27 & Friuli-Venezia Giulia & Italy     & Dialogue \\
28 & Liguria              & Italy      & Intimacy \\
29 & Piedmont             & Italy      & Philosophy \\
30 & Sardinia             & Italy      & Stubbornness \\
31 & Sicily               & Italy      & Resurrection \\
32 & Tuscany              & Italy      & Art \\
33 & Veneto               & Italy      & Commerce \\
34 & Douro                & Portugal   & Endurance \\
35 & Gori\v{s}ka Brda     & Slovenia   & Fortune \\
36 & Catalonia            & Spain      & Identity \\
37 & Galicia              & Spain      & Longing (Morri\~na) \\
38 & Ribera del Duero     & Spain      & Severity \\
39 & Rioja                & Spain      & Patience \\
40 & Mendoza              & Argentina  & Reinvention \\
41 & Patagonia            & Argentina  & Extremity \\
42 & Barossa Valley       & Australia  & Fortitude \\
43 & Hunter Valley        & Australia  & Defiance \\
44 & Margaret River       & Australia  & Composure \\
45 & Maipo Valley         & Chile      & Pride \\
46 & Central Otago        & NZ         & Adventure \\
47 & Hawke's Bay          & NZ         & Confidence \\
48 & Marlborough          & NZ         & Assertion \\
49 & Stellenbosch         & S.\ Africa & Aspiration \\
50 & Swartland            & S.\ Africa & Rebellion \\
51 & Columbia Valley      & USA        & Determination \\
52 & Finger Lakes         & USA        & Conviction \\
53 & Napa Valley          & USA        & Ambition \\
54 & Paso Robles          & USA        & Independence \\
55 & Santa Barbara        & USA        & Serendipity \\
56 & Santa Cruz Mountains & USA        & Obsession \\
57 & Sonoma               & USA        & Authenticity \\
58 & Walla Walla          & USA        & Community \\
59 & Willamette Valley    & USA        & Idealism \\
\bottomrule
\end{tabular}
\end{center}

\section*{D.2 Embedding Model Specification}
The embedding model used throughout is OpenAI's \texttt{text-embedding-3-small}. Specifications relevant to the Region Resonances deployment:

\begin{center}
\begin{tabular}{@{}lp{8.5cm}@{}}
\toprule
Property & Value \\
\midrule
Model identifier        & \texttt{text-embedding-3-small} \\
Output dimension        & $d = 1536$ \\
Output length           & Unit-normalised; $\norm{E(s)}_2 = 1$ for every input $s$ \\
Input language          & English (multilingual capability not used for matching) \\
Maximum input tokens    & 8{,}191 (well above any region passage length) \\
Determinism             & Same input always produces the same output \\
Treatment in pipeline   & Fixed, opaque function; weights and architecture not introspected \\
\bottomrule
\end{tabular}
\end{center}

\section*{D.3 Calibration Test Set and Operating Range}
The display transformation constants $s_{\min}$ and $s_{\max}$ were calibrated against a test set of 60 queries spanning four query types.

\begin{center}
\begin{tabular}{@{}lrp{8.0cm}@{}}
\toprule
Query type & $n$ & Examples \\
\midrule
Single abstract noun        & 18 & resilience, devotion, melancholy, pleasure \\
Compound abstract phrase    & 16 & old soul, quiet ambition, inherited burden \\
Sensory / atmospheric phrase & 14 & the feeling of late summer, dust on stone \\
Short prose fragment         & 12 & wines that ask you to wait \\
\midrule
Total                        & 60 & --- \\
\bottomrule
\end{tabular}
\end{center}

For each query, the full 59-element similarity vector was computed and the empirical distribution of cosine values recorded. Across all 60 queries, observed similarities fell almost entirely within $[0.15, 0.65]$. Distribution summary:

\begin{center}
\begin{tabular}{@{}lr@{}}
\toprule
Statistic & Value \\
\midrule
Min observed cosine (across all queries, all regions) & 0.08 \\
1st percentile                                          & 0.16 \\
Mean across all (query, region) pairs                  & 0.41 \\
99th percentile                                         & 0.66 \\
Max observed cosine (excluding paraphrase queries)     & 0.78 \\
Max observed cosine (paraphrase queries, $n = 4$)      & 0.91 \\
\bottomrule
\end{tabular}
\end{center}

Pathological inputs (random byte sequences, single-character queries, queries in languages other than English) routinely produce cosine values below 0.15. These are clipped to the lower display boundary by the transformation in Section D.4.

\section*{D.4 Display Transformation Constants}
The clip-and-rescale transformation is parameterised by two constants:
\[
\tilde{s}_i = \mathrm{clip}\!\left(\frac{s_i - s_{\min}}{s_{\max} - s_{\min}}, 0, 1\right) \times 100.
\]

\begin{center}
\begin{tabular}{@{}lr@{}}
\toprule
Constant & Value \\
\midrule
$s_{\min}$ (lower clip boundary)        & 0.15 \\
$s_{\max}$ (upper clip boundary)        & 0.70 \\
Display range (clipping aside)          & $[0, 100]$ \\
Cosine difference per display unit      & $0.0055$ \\
Display value for chance similarity ($s = 0.30$) & 27 \\
Display value for typical strong match ($s = 0.55$) & 73 \\
\bottomrule
\end{tabular}
\end{center}

The $s_{\min}$ value is set just above the chance baseline observed for unrelated queries; values below this are clipped to 0 to signal that the query has fallen outside the matcher's effective scope. The $s_{\max}$ value is set so that paraphrase queries (which produce cosine values above 0.70) saturate the display at 100; this prevents users from inferring fine distinctions among very-strong matches that the underlying cosine cannot reliably discriminate.

\section*{D.5 Distribution of Top-$k$ Cosine Similarities}
For the calibration test set, the distribution of top-1 and top-5 cosine similarities is summarised below. The top-1 distribution is the cosine of the best match for each query; the top-5 distribution is the cosine of the fifth-best match.

\begin{center}
\begin{tabular}{@{}lrr@{}}
\toprule
Statistic & Top-1 & Top-5 \\
\midrule
Min     & 0.42 & 0.32 \\
25th \%ile & 0.51 & 0.41 \\
Median  & 0.58 & 0.46 \\
Mean    & 0.59 & 0.47 \\
75th \%ile & 0.65 & 0.52 \\
Max     & 0.78 & 0.60 \\
\bottomrule
\end{tabular}
\end{center}

The typical top-1 to top-5 spread for a content-bearing query is approximately 0.10 in cosine units --- wide enough that the top match is consistently distinguishable from the fifth, while the fifth is consistently distinguishable from the next-best (sixth) region. After display transformation, top-1 scores range from approximately 50 to 100 with a median of 78; top-5 scores range from 30 to 80 with a median of 56. The display transformation preserves this rank-discriminating structure exactly.

\section*{References}

\begin{itemize}[leftmargin=2em,itemsep=0.2em]
\item OpenAI (2024). New embedding models and API updates. Technical announcement covering \texttt{text-embedding-3-small} and \texttt{text-embedding-3-large}.
\item Reimers, N. \& Gurevych, I. (2019). Sentence-BERT: Sentence embeddings using Siamese BERT-networks. In \emph{Proceedings of EMNLP-IJCNLP}, 3982--3992.
\end{itemize}

\end{document}
